% Part: proof-theory
% Chapter: sequent-calculus
% Section: quantifiers

\documentclass[../../../include/open-logic-section]{subfiles}

\begin{document}

\olfileid{pt}{seq}{qua}

\olsection{Bukti Reguler lan Substitusi}

Aturan kuantor ndadekake ana kerumitan amarga ``syarat
eigenvariabel'' kanggo aturan \LeftR{\lexists} lan~\RightR{\lforall}.
Umume kerumitan iki bisa ditangani kanthi njamin
!!{constant}~$a$ ing inferensi kasebut ora digunakake maneh.
Kita ngenalake dhefinisi iki.

\begin{defn}
!!^a{proof}~$\pi$ ing \Log{G1c} utawa \Log{G3c} diarani
\emph{reguler} yen
\begin{enumerate}
  \item ora ana eigenvariabel~$a$ sing digunakake minangka
  eigenvariabel luwih saka siji inferensi
  \LeftR{\lexists} utawa~\RightR{\lforall}, lan
  \item yen $a$ minangka eigenvariabel inferensi
  \LeftR{\lexists} utawa~\RightR{\lforall}, variabel iki
  mung ana ing ndhuwur inferensi kasebut ing~$\pi$.
\end{enumerate}
\end{defn}

\begin{lem}\ollabel{lem:subst-regular}
Yen $\pi$ reguler, dipungkasi ing $\Gamma \Sequent \Delta$,
$c$ minangka !!a{constant} sing ora digunakake minangka
eigenvariabel ing~$\pi$, lan $t$~minangka terma katutup sing ora
ngemot eigenvariabel saka~$\pi$, mula $\Subst{\pi}{t}{c}$
sing dijupuk saka~$\pi$ kanthi ngganti saben kedadeyan~$c$
nganggo~$t$ tanpa panangkepan variabel minangka !!{proof} reguler.
Sadurunge substitusi, pilih variabel kaiket sing seger
tumrap terma pangganti. Sekuen pungkasan
$\Subst{\pi}{t}{c}$ yaiku $\Subst{\Gamma}{t}{c} \Sequent
\Subst{\Delta}{t}{c}$, kanthi $\Subst{\Gamma}{t}{c} =
\Setabs{!A(t)}{!A(c) \in \Gamma}$. Notasi pangumpul iki
diwaca minangka multiset: saben kedadeyan asal menehi
siji kedadeyan sawise substitusi.
\end{lem}

\begin{proof}
  Kanthi induksi tumrap $\pheight{\pi}$.
  Yen $\pheight{\pi} = 0$, bukti mung kasusun saka aksioma.
  Mula $\Subst{\pi}{t}{c}$ uga aksioma, amarga saben
  !!{formula}~$!A(c)$ ing aksioma kasebut dadi $!A(t)$.

Yen $\pheight{\pi} > 0$, kita misahake kasus
  miturut inferensi pungkasan ing~$\pi$. Kasus aturan
  struktural (ing \Log{G1c}) lan aturan logis kanggo
  panyambung proposisional gampang lan ditinggal
  minangka latihan. Kita ngrembug kasus nalika
  $\pi$ dipungkasi nganggo inferensi kuantor.
  \begin{enumerate}
    \item Yen inferensi pungkasan yaiku $\RightR{\lforall}$,
    $\pi$ duwe wujud
    \[
    \AxiomC{}
    \RightLabel{$\pi'$}
    \Deduce$\Gamma(c) \fCenter \Delta'(c), !A(a,c)$
    \RightLabel{\RightR{\lforall}}
    \UnaryInf$\Gamma(c) \fCenter \Delta'(c), \lforall[x][!A(x,c)]$
    \DisplayProof
    \]
    Miturut hipotesis induksi, $\Subst{\pi'}{t}{c}$ minangka
    !!{proof} reguler kanggo $\Gamma(t) \fCenter \Delta'(t), !A(a,t)$.
    Amarga $a$ minangka eigenvariabel saka~$\pi$,
    $a$~ora ana ing~$t$. Mula inferensi pungkasan
    ing $\Subst{\pi}{t}{c}$,
    \[
    \AxiomC{}
    \RightLabel{$\Subst{\pi'}{t}{c}$}
    \Deduce$\Gamma(t) \fCenter \Delta'(t), !A(a,t)$
    \RightLabel{\RightR{\lforall}}
    \UnaryInf$\Gamma(t) \fCenter \Delta'(t), \lforall[x][!A(x,t)]$
    \DisplayProof
    \]
    bener: kesimpulane ora ngemot~$a$.
    \item Yen inferensi pungkasan yaiku $\LeftR{\lforall}$
    ing \Log{G3c}, $\pi$ duwe wujud
    \[
    \AxiomC{}
    \RightLabel{$\pi'$}
    \Deduce$!A(s(c),c), \lforall[x][!A(x,c)], \Gamma(c) \fCenter \Delta'(c)$
    \RightLabel{\LeftR{\lforall}}
    \UnaryInf$\lforall[x][!A(x,c)], \Gamma(c) \fCenter \Delta'(c)$
    \DisplayProof
    \]
    Miturut hipotesis induksi, $\Subst{\pi'}{t}{c}$ minangka
    !!{proof} reguler kanggo $!A(s(t),t), \lforall[x][!A(x,t)],
    \Gamma(t) \fCenter \Delta'(t)$. Inferensi pungkasan
    ing $\Subst{\pi}{t}{c}$,
    \[
    \AxiomC{}
    \RightLabel{$\Subst{\pi'}{t}{c}$}
    \Deduce$!A(s(t),t), \lforall[x][!A(x,t)], \Gamma(t) \fCenter \Delta'(t)$
    \RightLabel{\LeftR{\lforall}}
    \UnaryInf$\lforall[x][!A(x,t)], \Gamma(t) \fCenter \Delta'(t)$
    \DisplayProof
    \]
    bener. Kasus \LeftR{\lforall} ing~\Log{G1c} padha,
    nanging rada luwih prasaja.
    \item Inferensi pungkasan yaiku \RightR{\lexists}
    utawa \LeftR{\lexists}: minangka latihan.
  \end{enumerate}
  Ing saben kasus, kita nambahake sekuen ing pungkasan
  !!{proof} reguler (utawa rong !!{proof}s reguler
  kanggo aturan rong premis). Sekuen anyar beda saka
  sekuen pungkasan $\pi$ mung amarga $c$ wis diganti
  nganggo terma~$t$. Kajaba iku, eigenvariabel saka
  $\pi$ tetep: amarga $t$ ora ngemot eigenvariabel
  saka~$\pi$, eigenvariabel $\Subst{\pi}{t}{c}$
  padha karo sadurunge. Sekuen pungkasan anyar ora
  ngemot eigenvariabel $\Subst{\pi}{t}{c}$.
  Mula $\Subst{\pi}{t}{c}$ reguler.
\end{proof}

\begin{prop}
Yen $\pi$ minangka !!a{proof} ing \Log{G1c} utawa
\Log{G3c}, mula ana !!{proof} reguler~$\pi'$
kanthi struktur sing padha (mligine !!{height}
sing padha), sing beda saka $\pi$ mung amarga
eigenvariabele diganti.
\end{prop}

\begin{proof}
Mung kanggo ancas bukti iki, inferensi eigenvariabel
ing~$\pi$ diarani \emph{resik} yen eigenvariabele
ora dadi eigenvariabel inferensi liyane lan ora
ana ing panggonan liyane ing~$\pi$ kajaba ing
sub-!!{proof} langsung saka inferensi kasebut;
yen ora, inferensi diarani \emph{reged}.
Upamane $n$ minangka cacah inferensi eigenvariabel
reged ing~$\pi$. Kita nuduhake kanthi induksi
kuwat tumrap~$n$ yen $\pi$~bisa diowahi dadi
!!{proof} reguler~$\pi'$ sing dibutuhake.
Yen $n=0$, kabeh inferensi eigenvariabel ing~$\pi$
resik, mula $\pi$ reguler lan ora ana sing perlu
dibuktekake. Yen ora, pilih inferensi eigenvariabel
reged sing paling ndhuwur, upamane \RightR{\lforall}.
Sub-!!{proof}~$\pi_1$ saka~$\pi$ sing dipungkasi
nganggo inferensi iki duwe wujud
\[
    \AxiomC{}
    \RightLabel{$\pi_2$}
    \Deduce$\Gamma \fCenter \Delta, !A(c)$
    \RightLabel{\RightR{\lforall}}
    \UnaryInf$\Gamma \fCenter \Delta, \lforall[x][!A(x)]$
    \DisplayProof
    \]
Gatekna yen amarga $c$ minangka eigenvariabel,
konstanta iki ora ana ing~$\Gamma$, $\Delta$,
utawa $\lforall[x][!A(x)]$. Bukti $\pi_2$ reguler,
amarga kabeh inferensi eigenvariabele resik lan
ora ana sing nggunakake konstanta sing dipilih
minangka eigenvariabel. Pilih $a$ minangka
!!a{constant} sing ora ana ing~$\pi$.
Miturut \olref{lem:subst-regular},
$\Subst{\pi_2}{a}{c}$ minangka bukti reguler
kanggo $\Gamma \fCenter \Delta, !A(a)$,
lan kita bisa ngetrapake~$\RightR{\lforall}$.
Yen $\pi_1$ ing~$\pi$ diganti nganggo bukti
$\pi_1'$,
\[
    \AxiomC{}
    \RightLabel{$\Subst{\pi_2}{a}{c}$}
    \Deduce$\Gamma \fCenter \Delta, !A(a)$
    \RightLabel{\RightR{\lforall}}
    \UnaryInf$\Gamma \fCenter \Delta, \lforall[x][!A(x)]$
    \DisplayProof
    \]
kita wis ngganti inferensi eigenvariabel reged
dadi inferensi resik, lan prabedane mung
panggantian eigenvariabel $c$ nganggo~$a$.
Bukti asile duwe inferensi eigenvariabel reged
sing cacahe kurang saka $n$, lan asil sing
dibutuhake ditampa saka hipotesis induksi.
\end{proof}

Kita ora bakal nyebut proposisi sing lagi
dibuktekake iki kanthi eksplisit, nanging bakal
nganggep kabeh !!{proof}s sing dirembug reguler---
yen durung reguler, eigenvariabele mesthi bisa
diganti supaya dadi reguler.

\end{document}
